\setchapterpreamble[u]{\margintoc}
\chapter{Conclusions}
\labch{Con}

This thesis explored the application of ST technologies to characterise the tumour microenvironment through a complete computational workflow including QC, cell-type annotation, and neighbourhood analysis. Despite the limitations of targeted ST panels, the results demonstrated that biologically meaningful spatial patterns can still be identified when robust analytical strategies are applied.

The quality-control stage allowed the identification and filtering of low-quality cells and technical artefacts while preserving biologically meaningful spatial information. Different metrics were evaluated to establish a robust and reproducible filtering strategy adapted to Xenium data. This step was especially important because the quality of downstream analyses strongly depended on maintaining a balance between removing noisy cells and preserving tissue structure.

Cell-type annotation represented another important part of the work. The results showed that unsupervised clustering alone was insufficient to accurately define all cellular populations due to overlap between transcriptomic profiles in low-dimensional datasets. For this reason, multiple complementary strategies were combined, including marker-gene exploration, enrichment analyses, and reference-based annotation methods. This integrated approach improved the biological interpretation of the data and allowed the identification of relevant stromal, epithelial, and immune populations within the samples.

One of the most important results of the project emerged from the neighbourhood analysis. The study of spatial cellular interactions revealed non-random tissue organisation patterns and suggested potential biological mechanisms associated with treatment response. In particular, the analyses showed spatial enrichment of fibroblast populations surrounding or interacting with epithelial-enriched relevant cellular communities in samples with poorer therapeutic response. Based on these observations, the results obtained in this work further supported an existing hypothesis suggesting that fibroblasts could contribute to immune exclusion and could represent one mechanism of resistance to ICI in MIBC. However, larger validation cohorts, direct quantification of tumour–immune cell distances, cancer associated fibroblasts phenotyping, integration with additional biomarkers, and functional studies would be required to demonstrate a causal role of fibroblasts in ICI resistance.

Nevertheless, these findings must be interpreted cautiously. Although the neighbourhood analyses revealed biologically meaningful associations, additional statistical validation and larger cohorts would be necessary to confirm any causal relationship between fibroblast spatial organisation and treatment resistance. Furthermore, some samples included in the dataset correspond to patients with complete pathological response or cases without detectable residual tumour despite their original clinical classification. Therefore, the observed spatial patterns should not automatically be interpreted as direct evidence of active tumour presence.

Although this work builds upon existing methodologies previously described in the literature, its innovative aspect lies in applying them to address a clinically relevant question arising from the DUTRENEO trial, one of the first prospective studies evaluating biomarker-guided immunotherapy in muscle-invasive bladder cancer (MIBC). While biomarkers such as the Tumour Inflammation Score (TIS) have shown promise in retrospective studies, their prospective predictive value remains limited. By using spatial transcriptomics to investigate the tumour microenvironment and identify mechanisms such as immune exclusion that may underlie treatment resistance, this work contributes to the development of more accurate patient stratification strategies and advances the implementation of personalized medicine in MIBC. Another innovative aspect resides in the integration and standardisation of a complete ST analysis workflow focused specifically on Xenium data. The novelty is therefore not centred on the invention of entirely new algorithms, but rather on the practical adaptation, evaluation, and harmonisation of existing approaches for a rapidly evolving technological field where standardised analytical practices are still under development. In this sense, the project also contributes to improving the reproducibility and interpretability of ST analyses.  

Regarding sustainability, this project relied entirely on computational analyses of existing datasets, reducing the need for additional laboratory experimentation and biological material consumption. Nevertheless, large-scale computational analyses still involve significant computational resources and energy usage, making efficient pipeline design increasingly important.

From an ethical and social perspective, ST has strong potential for advancing precision medicine and improving our understanding of complex diseases such as cancer. Since the gene panels used in both Visium and Xenium technologies only capture a limited subset of the genetic information present in the cells, the data generated is not sufficient to identify individual patients. Therefore, the genetic information used in this thesis does not represent a major security risk in terms of patient identification. Nevertheless, patient consent and responsible data management must always remain a priority. The patients participating in the DUTRENEO clinical trial were informed about the use of their data, provided consent for its use in research, and all personal information was anonymised before analysis. Additionally, patients who were considered unlikely to benefit from ICI therapy, based on their low TIS ("cold" tumours), were assigned to the standard-of-care treatment rather than being allocated to a control group with limited prospects of therapeutic benefit. This study design avoided exposing these patients to potentially ineffective therapy and ensured they were not deprived of established treatment options.

Additionally, the presence of samples showing complete pathological response despite their original tumour classification, or samples classified as non-responsive without detectable residual tumour, should not be interpreted as false positives or false negatives. Instead, these situations reflect the complexity and limited explainability of the available clinical and biological data, as well as the inherent challenges in interpreting treatment response within heterogeneous tumour microenvironments.

Finally, diversity and representation remain important challenges in ST research. Although the dataset used in this thesis was multicentric, all samples originated from Spanish hospitals. This introduces potential demographic and biological biases, as the cohort may not fully represent broader population variability. Expanding cohort diversity in future studies will therefore be essential to ensure that the resulting biological insights and potential clinical applications are broadly representative and equitable.

Overall, this work demonstrates that robust and interpretable biological insights can be extracted from targeted ST datasets through careful methodological design and critical interpretation, providing a reproducible basis for future studies in this rapidly developing field. Nevertheless, the results support the hypothesis that stromal-mediated immune exclusion may represent one mechanism underlying the limited predictive value of TIS observed in the DUTRENEO clinical trial.


\section{Code Availability}
All the code used in this thesis can be accessed in this GitHub repository:
\url{https://github.com/SofiaDiCapua/Spatial-Transcriptomics-Masther-Thesis}

Artificial Intelligence (AI) tools were used during the development of this Master’s Thesis and in the preparation of this document. ChatGPT was used to review and correct the final drafts of the different sections of the thesis; it was not used to generate the text from scratch, and all suggested corrections were carefully reviewed before being incorporated. Additionally, Claude was used to assist in the development of the more complex code components and in debugging particularly challenging errors.

\section{Acknowledgements}
First and foremost, I would like to express my immense gratitude to all the people who supported me throughout the long process of developing this master's thesis project. I would especially like to thank my partner, family, and close friends for their endless support and encouragement.

I would also like to thank the Cancer Immunogenomics group at the Josep Carreras Leukaemia Research Institute for welcoming me into the group and for all the valuable insights gained through the many discussions surrounding this project. A very special thanks goes to Marc Escobosa, Sergi Cervilla, and Ivo Cruz for answering my endless questions, as well as for their constant support and supervision throughout this work.

Special thanks also go to Daniela Grases and Elena Pérez for their exceptional work in the wet lab with the samples analysed in this thesis, and to Eduard Porta for his continuous supervision, guidance, and invaluable advice throughout the entire process. Finally, I would also like to thank Mustafa Sibai for his patience and for all the time and dedication he devoted to helping me during this project.

%%%%%%%%%%%%%%%%%%%%%%%%%%%%%%%%%%%%%%%%%%%%%%%%%%%%%%%%%%%%%%%%%%%%%%%%%%%%%%%%%%%%%%%%%%%%%%%

% Sidenotes are a distinctive feature of all 1.5-column-layout books. 
% Indeed, having wide margins means that some material can be displayed 
% there. We use margins for all kind of stuff: sidenotes, marginnotes, 
% small tables of contents, citations, and, why not?, special boxes and 
% environments.

% \section{Sidenotes}

% Sidenotes are like footnotes, except that they go in the margin, where 
% they are more readable. To insert a sidenote, just use the command 
% \Command{sidenote\{Text of the note\}}. You can specify a 
% mark\sidenote[O]{This sidenote has a special mark, a big O!} with \\ 
% \Command{sidenote[mark]\{Text\}}, but you can also specify an offset, 
% which moves the sidenote upwards or downwards, so that the full syntax is:

% \begin{lstlisting}[style=kaolstplain]
% \sidenote[mark][offset]{Text}
% \end{lstlisting}

% If you use an offset, you always have to add the brackets for the mark, 
% but they can be empty.\sidenote{If you want to know more about the usage 
% of the \Command{sidenote} command, read the documentation of the 
% \Package{sidenotes} package.}

% In \Class{kaobook} we copied a feature from the \Package{snotez} 
% package: the possibility to specify a multiple of \Command{baselineskip} 
% as an offset. For example, if you want to enter a sidenote with the 
% normal mark and move it upwards one line, type:

% \begin{lstlisting}[style=kaolstplain]
% \sidenote[][*-1]{Text of the sidenote.}
% \end{lstlisting}

% As we said, sidenotes are handled through the \Package{sidenotes} 
% package, which in turn relies on the \Package{marginnote} package.

% \section{Marginnotes}

% This command is very similar to the previous one. You can create a 
% marginnote with \Command{marginnote[offset]\{Text\}}, where the offset 
% argument can be left out, or it can be a multiple of 
% \Command{baselineskip},\marginnote[-1cm]{While the command for margin 
% notes comes from the \Package{marginnote} package, it has been redefined 
% in order to change the position of the optional offset argument, which 
% now precedes the text of the note, whereas in the original version it 
% was at the end. We have also added the possibility to use a multiple of 
% \Command{baselineskip} as offset. These things were made only to make 
% everything more consistent, so that you have to remember less things!} 
% \eg

% \begin{lstlisting}[style=kaolstplain]
% \marginnote[-12pt]{Text} or \marginnote[*-3]{Text}
% \end{lstlisting}

% \begin{kaobox}[frametitle=To Do]
% A small thing that needs to be done is to renew the \Command{sidenote} 
% command so that it takes only one optional argument, the offset. The 
% special mark argument can go somewhere else. In other words, we want the 
% syntax of \Command{sidenote} to resemble that of \Command{marginnote}.
% \end{kaobox}

% We load the packages \Package{marginnote}, \Package{marginfix} and 
% \Package{placeins}. Since \Package{sidenotes} uses \Package{marginnote}, 
% what we said for marginnotes is also valid for sidenotes. Side- and 
% margin- notes are shifted slightly upwards 
% (\Command{renewcommand\{\textbackslash marginnotevadjust\}\{3pt\}}) in 
% order to align them to the bottom of the line of text where the note is 
% issued. Importantly, both sidenotes and marginnotes are defined as 
% floating if the optional argument (\ie the vertical offset) is left 
% blank, but if the offset is specified they are not floating. Recall that 
% floats cannot be nested, so in some rare cases you may encounter errors 
% about lost floats; in those cases, remember that sidenotes and 
% marginnotes are floats. To solve the problem, it may be possible to 
% transform them into non-floating elements by specifying an offset of 
% 0pt.

% \section{Footnotes}

% Even though they are not displayed in the margin, we will discuss about 
% footnotes here, since sidenotes are mainly intended to be a replacement 
% of them. Footnotes force the reader to constantly move from one area of 
% the page to the other. Arguably, marginnotes solve this issue, so you 
% should not use footnotes. Nevertheless, for completeness, we have left 
% the standard command \Command{footnote}, just in case you want to put a 
% footnote once in a while.\footnote{And this is how they look like. 
% Notice that in the PDF file there is a back reference to the text; 
% pretty cool, uh?}

% \section{Margintoc}

% Since we are talking about margins, we introduce here the 
% \Command{margintoc} command, which allows one to put small table of 
% contents in the margin. Like other commands we have discussed, 
% \Command{margintoc} accepts a parameter for the vertical offset, like 
% so: \Command{margintoc[offset]}.

% The command can be used in any point of the document, but we think it 
% makes sense to use it just at the beginning of chapters or parts. In 
% this document I make use of a \KOMAScript\xspace feature and put it in 
% the chapter preamble, with the following code:

% \marginnote{The font used in the margintoc is the same as the one for 
% 	the chapter entries in the main table of contents at the beginning 
% 	of the document.}

% \begin{lstlisting}[style=kaolstplain]
% \setchapterpreamble[u]{\margintoc}
% \chapter{Chapter title}
% \end{lstlisting}

% As the space in the margin is a valuable resource, there is the 
% possibility to print a shorter version of the title in the margin toc. 
% Thus, there are in total three possible versions for the title of a 
% section (or subsection): the one for the main text, the one for the main 
% table of contents, and the one for the margintoc. These versions can be 
% specified at the same time when the section is created in the source 
% \TeX file:
% \begin{lstlisting}[style=kaolstplain]
% \section[alternative-title-for-toc]{title-as-written-in-text}[alternative-title-for-margintoc]
% \end{lstlisting}

% By default, the margintoc includes sections and subsections.
% If you only want to show sections, add
% \begin{lstlisting}[style=kaolstplain]
% \setcounter{margintocdepth}{\sectiontocdepth}
% \end{lstlisting}
% somewhere in your preamble.

% \section{Marginlisting}

% On some occasions it may happen that you have a very short piece of code 
% that doesn't look good in the body of the text because it breaks the 
% flow of narration: for that occasions, you can use a 
% \Environment{marginlisting}. The support for this feature is still 
% limited, especially for the captions, but you can try the following 
% code:

% \begin{marginlisting}[-1.35cm]
% 	\caption{An example of a margin listing.}
% 	\vspace{0.6cm}
% 	\begin{lstlisting}[language=Python,style=kaolstplain]
% print("Hello World!")
% 	\end{lstlisting}
% \end{marginlisting}

% \begin{verbatim}
% \begin{marginlisting}[-0.5cm]
% 	\caption{My caption}
% 	\vspace{0.2cm}
% 	\begin{lstlisting}[language=Python,style=kaolstplain]
% 	... code ...
% 	\end{lstlisting}
% \end{marginlisting}
% \end{verbatim}

% Unfortunately, the space between the caption and the listing must be 
% adjusted manually; if you find a better way, please let me know.

% Not only textual stuff can be displayed in the margin, but also figures. 
% Those will be the focus of the next chapter.
